\section{Large fibers of the totient function}\label{sec:totients}

We finish by deriving \Cref{thm:totient} from \Cref{thm:smooth}.
This product-and-pigeonhole passage is the classical connection between
smooth shifted primes and large totient fibers; see
\cite[Theorem~B and its proof]{Pomerance1980}.
We first record finiteness of each fiber and the elementary upper
bound mentioned in the Introduction.

\begin{lemma}\label{lem:finite-totient-fibers}
For each positive integer $n$, the set of positive integers $m$ with
$\varphi(m)=n$ is finite. For every $\eta>0$,
\[
 g(n)\ll_\eta n^{1+\eta}.
\]
\end{lemma}

\begin{proof}
If $p^a\parallel m$, then $p^{a-1}(p-1)$ divides $\varphi(m)$.
For a fixed value $n$, this restricts $p$ to the finite set with
$p-1\mid n$, and bounds $a$ by $p^{a-1}\mid n$. Thus each fiber
is finite.

For the quantitative bound, fix $0<\theta<1$. For all sufficiently
large primes $p$, one has $1-1/p\ge p^{-\theta}$. The finitely many
remaining primes contribute a fixed positive constant $c_\theta$.
Consequently, for every $m$,
\[
 \varphi(m)=m\prod_{p\mid m}\left(1-\frac1p\right)
 \ge c_\theta m\prod_{p\mid m}p^{-\theta}
 \ge c_\theta m^{1-\theta}.
\]
Every preimage of $n$ is therefore at most
$(n/c_\theta)^{1/(1-\theta)}$. Choose $\theta$ so that
$1/(1-\theta)\le1+\eta$, and count the possible positive integers.
\end{proof}

\begin{proof}[Proof of \Cref{thm:totient}]
It suffices to consider $0<\eps<1$. Choose $0<\delta<1/4$ with
$4\delta<\eps$. By \Cref{thm:smooth}, for every sufficiently large
$X$ there are more than $X^{1-\delta}$ primes $p\le X$ such that
$p-1$ is $X^\delta$-smooth. Indeed, take $x=X/5$ in that theorem;
its count $x^{1-o(1)}$ exceeds $X^{1-\delta}$ eventually, and
$x^\delta\le X^\delta$.

Let $\mathcal P$ be this set of primes, let $M=|\mathcal P|$, and
put $k=\lfloor X^{2\delta}\rfloor$. Since $3\delta<1$, we have
$k\le M$ for sufficiently large $X$. Different $k$-element subsets
of $\mathcal P$ have different products by unique factorization.
There are at least
\begin{equation}\label{eq:many-prime-products}
 \binom Mk\ge\left(\frac Mk\right)^k
 >X^{(1-3\delta)k}
\end{equation}
such products. For completeness, the binomial inequality follows
by writing $\binom Mk=\prod_{j=0}^{k-1}(M-j)/(k-j)$ and observing
that each factor is at least $M/k$.

For a squarefree product $m=\prod_{p\in\mathcal Q}p$ with
$|\mathcal Q|=k$, multiplicativity of the totient gives
\[
 \varphi(m)=\prod_{p\in\mathcal Q}(p-1).
\]
This value is $X^\delta$-smooth and at most $X^k$. Every prime
exponent in such a value is at most $k\log X/\log2$, and there
are at most $X^\delta$ primes available. Hence the total number
of possible values is at most
\begin{equation}\label{eq:few-totient-values}
 \left(1+\frac{k\log X}{\log2}\right)^{X^\delta}
 =X^{o(k)}.
\end{equation}
The last equality means that the logarithm of its left side,
divided by $k\log X$, tends to zero: its numerator is
$O_\delta(X^\delta\log X)$, whereas $k\asymp X^{2\delta}$.

By \Cref{eq:many-prime-products,eq:few-totient-values}, some
$n\le X^k$ therefore satisfies
\[
 g(n)\ge X^{(1-3\delta-o(1))k}>n^{1-\eps}
\]
for sufficiently large $X$. The strict final inequality follows
from $3\delta+o(1)<\eps$ and $1-\eps>0$.

These lower bounds for $g(n)$ tend to infinity with $X$.
By \Cref{lem:finite-totient-fibers}, the resulting values $n$ cannot
belong to a fixed finite set. They are therefore unbounded, which
proves the required assertion for infinitely many $n$. If
$\eps\ge1$, the assertion follows from any smaller positive
choice of $\eps$.
\end{proof}
